When we compute LO DIS from \cref{fig:LO} we say it is easy to identify the respective coefficient function, \cref{eq:LObarCq,eq:LOC}.
However, what exactly do we mean with the \enquote{$\forall q$}?
Which quarks are we exactly talking about?
This is directly linked to the flavor space to which we hint when using bold font.
However, unlike the problem we discuss in \cref{sec:svimpl} about different bases (flavor basis vs.\ evolution basis), here we are more directly concerned what the flavor space is in the first place.

\subsection{Flavor space}

We define the flavor space $\mathcal F$ as the vector space spanned by all quark flavors and the gluon.
We consider the naive flavor basis as our canonical basis and write
\begin{equation}
    \mathcal F = \mathcal F_{fl} = \mathrm{span}(g, u, \bar u, d, \bar d, s, \bar s, c, \bar c, b, \bar b, t, \bar t)\,.
\end{equation}
We have a 13 dimensional vector space, $\dim{\mathcal F} = 13$\footnote{one can add the photon and make it 14 dimensional~\cite{NNPDF:2024djq}}, consisting of the gluon, six quarks, and six anti-quarks.
Our bold vectors are an element in that space, e.g.\ $\vb C \in \mathcal F$, or more plainly,
\begin{equation}
    \vb C^T \sim (C_g, C_u, C_{\bar u}, C_d, C_{\bar d}, C_s, C_{\bar s}, C_c, C_{\bar c}, C_b, C_{\bar b}, C_t, C_{\bar t})
\end{equation}
and, as said in \cref{sec:svimpl}, any other basis works just as well.
The same applies also to the PDF $\vb f$, which is the other vector in \cref{eq:fact3}, and any operator, e.g.\ splitting functions $\vb P$, just lives in the product space, $\vb P \in \mathcal F \otimes \mathcal F$.

The question we need to address here is: which quarks exactly can participate in \cref{eq:LOxs,eq:LOC}?
We refer to this number as the number of active quarks $n_f$.
This dependency is completely hidden in \cref{eq:fact1}, but, as we demonstrate in the following, is a major concern for both coefficient functions and PDF.

The number of flavors $n_f$ is arbitrary per se and there is no absolutely correct answer, since there is no theory argument, which would fix this choice.
However, in practice for a given observable $\sigma(x,Q^2)$ we can consider, which options are phenomenologically reasonable and if some option would yield unstable results.

For the sake of the argument, let's consider a fully inclusive DIS cross section $\sigma(x,Q^2=\SI{25}{\GeV^2})$, so we probe the probe the proton at $\mu_F=Q=\SI{5}{\GeV}$.
We can now consider the different quark masses as a guideline: on the one side the masses of up, down, and strange are significantly below the non-perturbative regime of QCD $m_u,m_d,m_s \ll \LQCD \sim \SI{1}{\GeV}$ and we can safely assume we can find them in the proton.
On the other side the mass of the top is significantly above the probed scale, $m_t \gg Q$, thus it is very unlikely to find it in the proton and we can neglect it.
We can shut down such contributions by setting $f_t(\xi, Q^2) = 0 = f_{\bar t}(\xi, Q^2)$ and this way the top quark does not contribute to \cref{eq:fact1}.
In addition, we can also define $C_t(z, Q^2) = 0 = C_{\bar t}(z, Q^2)$ to explicitly say that we don't let the top contribute.

Now, we consider the two remaining quarks, charm and bottom, for which maybe there is no clear answer.
First, for the charm quark we have $m_c = \SI{1.5}{\GeV}$ and we can assume reasonably that it can contribute to DIS.
On the other side, for the bottom quark we have $m_b = \SI{4.5}{\GeV}$ and maybe we can both allow or deny it.
The specific choice on which quarks we allow to participate and in which way they do so is referred to as flavor number scheme.

\subsection{ZM-VFNS}

One very popular choice is the Zero-Mass Variable Flavor Number Scheme (ZM-VFNS), which is a telling name.
Formally, we can define it in the following way:
first, whenever the probed scale $Q$ crosses a quark mass that quark becomes active;
second; in coefficient functions all active quarks are considered massless and all non-active quarks are considered infinite massive.
The infinite mass ensures $f_H(\xi, Q^2) = 0 = f_{\bar H}(\xi, Q^2)$ and $C_H(z, Q^2) = 0 = C_{\bar H}(z, Q^2)$ for all non-active quarks $H$ and the zero mass ensures that all formulas in \cref{sec:nlo} hold.
If we denote the new dependency on the number of active flavors $n_f$ by a superscript $(n_f)$ we can define the ZM-VFNS in the following way
\begin{equation}
    \sigma^{\text{ZM-VFNS}}(x,Q^2) = \left\{ \begin{array}{ll}
        \sigma^\mathrm{(3)}(x,Q^2) &\text{if}~ Q^2 < m_c^2\\
        \sigma^\mathrm{(4)}(x,Q^2) &\text{if}~ m_c^2 \leq Q^2 < m_b^2\\
        \sigma^\mathrm{(5)}(x,Q^2) &\text{if}~ m_b^2 \leq Q^2 < m_t^2\\
        \sigma^\mathrm{(6)}(x,Q^2) &\text{if}~ Q^2 \geq m_t^2
    \end{array} \right. \label{eq:zmvfns}
\end{equation}
where
\begin{equation}
    \sigma^{(n_f)}(x,Q^2) = \left(\vb C^{T,(n_f)}(Q^2) \otimes \vb f^{(n_f)}(Q^2)\right)(x)
\end{equation}
i.e.\ each contribution has its own factorization formula and we have, e.g.,
\begin{equation}
    \begin{array}{l}
        C_H^{(3)}(z, Q^2) = 0 = C_{\bar H}^{(3)}(z, Q^2)\\
        f_H^{(3)}(\xi, Q^2) = 0 = f_{\bar H}^{(3)}(\xi, Q^2)
    \end{array}
    \quad H\in\{c,\bar c, b, \bar b, t, \bar t\}
\end{equation}
and similar for the others.
Specifically, we can improve \cref{eq:LOC} by writing
\begin{align}
    C_q^{(0),(3)}(z) &= e_q^2 \delta(1-z) & q\in\{c,\bar c, b, \bar b, t, \bar t\}\\
    C_H^{(0),(3)}(z) &= 0 & H\in\{c,\bar c, b, \bar b, t, \bar t\}\\
    C_g^{(0),(3)}(z) &= 0
\end{align}
and we finally give full details on the available flavor space $\mathcal F$.

Note that it does not make sense to consider less then three quarks, since $m_u,m_d,m_s \ll \LQCD \sim \SI{1}{\GeV}$ and pQCD only works for $Q > \LQCD$.
All quarks are considered with their actual mass for the decision if they are in or out, \cref{eq:zmvfns}, but then after that their mass is discarded when computing $\vb C^{(n_f)}(z,Q^2)$.
Thus, the VFNS indicates that the number of active quarks $n_f$ depends on the probed scale $Q$, i.e.\ it is variable, and the ZM explains how to compute diagrams.
However, once we have an actual scale in hand, $n_f$ is also fixed.

\subsection{FFNS}

For the number of active quarks we can make another simple choice: fix $n_f$ once and forever - the Fixed Flavor Number Scheme (FFNS).
Here, the probed scale $Q$ does not matter, but we use the same coefficient functions for any scale.
Actually, saying FFNS is not sufficient, we still have some choices and the most obvious is $n_f$.

However, we can also improve with respect to other parts: while it is significantly easier to compute diagrams with only massless quarks, we can compute them also with massive quarks.
We recall two important features of the propagator of a massive quark,
\begin{equation}
    D_q(p) \sim \frac{p_\mu\gamma^\mu + m\mathbbm{1}}{p^2 - m^2}\,;
\end{equation}
first, the mass appears explicitly in the numerator, which makes computing Dirac traces significantly more lengthy, and, second, the mass appears in the denominator, which shields the propagator from a collinear pole.
When we consider a finite mass we eventually have three different scales in our diagrams: the total energy of the colliding partons $\hat s$, the photon virtuality $Q^2$, and the quark mass $m^2$.
This means that unlike the massless case discussed so far, the coefficient functions now depend on two scales, for example the partonic momentum fraction $z = Q^2/(\hat s + Q^2)$ and $\xi_q = Q^2 / m^2$ or any other bijective combination\footnote{For example, for NC DIS on can show~\cite{Hekhorn:2019nlf} that $\chi = \frac{1-\sqrt{1-4m^2/s}}{1+\sqrt{1-4m^2/s}}$ and $\chi_q = \frac{\sqrt{1+4m^2/Q^2}-1}{\sqrt{1+4m^2/Q^2}+1}$ are good combinations.}.
This in turn implies that the available mathematical function space is much bigger, which makes the computation of these coefficient functions $\vb C(z,\xi_q,Q^2)$ significantly more complex.

However, we focus here on the new dependency on $\xi_q$, which captures the new features.
First, we consider the case $Q^2 \gg m^2$ or $\xi_q \to \infty$:

\subsection{FONLL}

\subsection{Heavy quark structure functions}
